% T27 — «Лес/природа»: зелёная гамма, листики-маркеры, кремовый фон.
% Math, класс 6, 6 задач (дроби и проценты).
\documentclass[a4paper,12pt]{article}

\usepackage{../../../../Lessons/_templates/preamble_worksheet}
\usepackage{../_shared/worksheet_check}

\usepackage{eso-pic}

% --- Палитра: лесной зелёный + кремовый ---
\definecolor{wcprimary}{HTML}{2F6E3F}
\definecolor{wcprimaryLight}{HTML}{B8D9B8}
\definecolor{wcprimaryBG}{HTML}{FBF6E7}
\definecolor{wcleaf}{HTML}{6CA86B}
\definecolor{wcbark}{HTML}{8B6A4A}

\renewcommand{\familydefault}{\sfdefault}

% --- Кремовый фон страницы ---
\AddToShipoutPictureBG{%
  \AtPageLowerLeft{%
    \color{wcprimaryBG}\rule{\paperwidth}{\paperheight}}}

% --- Заголовок с листиками ---
\renewcommand{\WorksheetTitle}[1]{%
  \par\noindent\begin{tikzpicture}
    % Зелёная полоса
    \fill[wcprimary] (0,0) rectangle (\linewidth, -16mm);
    % Листики справа (овальные)
    \begin{scope}[shift={([xshift=-22mm,yshift=-8mm]\linewidth,0)}]
      \fill[wcleaf, rotate=-25] (0,0) ellipse (5mm and 2.2mm);
      \fill[wcleaf!85, rotate=20]  (3mm,1mm) ellipse (4mm and 1.8mm);
      \fill[wcleaf!70, rotate=-5]  (6mm,-1mm) ellipse (4.5mm and 2mm);
    \end{scope}
    \node[anchor=west, text=white, font=\bfseries\Large\sffamily]
      at (4mm, -5mm) {ЛЕСНАЯ МАТЕМАТИКА};
    \node[anchor=west, text=wcprimaryLight, font=\bfseries\large\sffamily]
      at (4mm, -11mm) {#1};
  \end{tikzpicture}\par\vspace{3mm}}

% --- TaskBox: листик-маркер с номером ---
\renewcommand{\TaskBox}[2]{%
  \par\noindent\begin{tikzpicture}
    \node[draw=wcleaf, line width=0.9pt, rounded corners=5pt,
          fill=white, inner sep=3.5mm,
          text width=\dimexpr\linewidth-18mm\relax,
          anchor=north west] (B) at (10mm,0) {#2};
    % Листик
    \fill[wcleaf, rotate around={-20:(4mm,-4mm)}]
      (4mm,-4mm) ellipse (5mm and 3mm);
    \draw[wcprimary, line width=0.8pt, rotate around={-20:(4mm,-4mm)}]
      (-1mm,-4mm) -- (9mm,-4mm);
    \node[font=\bfseries\color{white}] at (4mm,-4mm) {#1};
  \end{tikzpicture}\par\vspace{3mm}}

\SetAnswerFieldSize{30mm}{8mm}

\begin{document}

\WorksheetTitle{Дроби и проценты в лесу}
{\small\itshape\color{wcprimary!85}Класс 6 \quad·\quad T27 \quad·\quad ID: T27-forest-2026-05-27}\par\vspace{2mm}
\FirstPageHeader

\TaskBox{1}{%
  В лесу растёт $120$ деревьев. $\dfrac{1}{4}$ из них~--- ели.
  Сколько в лесу елей?\par
  \vspace{1.5mm}\hfill\textbf{\color{wcprimary}Ответ:}~\AnswerField{1}%
}

\TaskBox{2}{%
  Белка собрала $80$ орехов. $25\%$ из них она спрятала в дупло.
  Сколько орехов она спрятала?\par
  \vspace{1.5mm}\hfill\textbf{\color{wcprimary}Ответ:}~\AnswerField{2}%
}

\TaskBox{3}{%
  Ёжик прошёл $\dfrac{2}{5}$ пути за час. Весь путь~--- $25$~км.
  Сколько километров он прошёл?\par
  \vspace{1.5mm}\hfill\textbf{\color{wcprimary}Ответ:}~\AnswerField{3}%
}

\TaskBox{4}{%
  В кленовом букете $40$ листьев. $30\%$ из них~--- красные, остальные~--- жёлтые.
  Сколько жёлтых листьев?\par
  \vspace{1.5mm}\hfill\textbf{\color{wcprimary}Ответ:}~\AnswerField{4}%
}

\TaskBox{5}{%
  Лесник посадил $\dfrac{3}{8}$ от $480$ саженцев в первый день.
  Сколько саженцев он посадил?\par
  \vspace{1.5mm}\hfill\textbf{\color{wcprimary}Ответ:}~\AnswerField{5}%
}

\TaskBox{6}{%
  В заповеднике $50$ лосей. После расселения их стало на $20\%$ больше.
  Сколько лосей стало в заповеднике?\par
  \vspace{1.5mm}\hfill\textbf{\color{wcprimary}Ответ:}~\AnswerField{6}%
}

\WorksheetQR{T27-forest/L1/v1}

\end{document}
